\documentclass[12pt]{article}
\usepackage[utf8]{inputenc}
\usepackage[T1]{fontenc}
\usepackage[brazil]{babel}
\usepackage[a4paper,margin=2.5cm]{geometry}
\usepackage{enumitem}
\setlist[itemize]{noitemsep,topsep=4pt}
\title{Sprint 3 - S3\_03\\Filas e Listas Encadeadas}
\author{Matheus Moreli Durão Silva\\Matrícula 12621CBS217\\Cibersegurança - Bacharelado - Noturno}
\date{2026}
\begin{document}
\maketitle

\section*{Entregável}
Implementação de uma fila e uma lista encadeada com exemplos. Este documento registra o objetivo do card, a solução desenvolvida e a forma de conferência. Os arquivos de código e as evidências citadas permanecem na mesma pasta deste documento.

\section*{O que foi estudado e aplicado}
\begin{itemize}
\item FIFO: inserir no fim e remover do início
\item Nó com valor e próximo
\item Inserção e remoção
\item Complexidade simplificada das operações
\end{itemize}

\section*{Como conferir}
Execute \texttt{python estruturas.py}. A fila remove primeiro o item inserido primeiro; a lista encadeada insere três nós, remove o nó intermediário e percorre os restantes. \texttt{diagrama\_estruturas.txt} representa as referências. Nesta implementação, inserir na fila é $O(1)$ amortizado, \texttt{pop(0)} é $O(n)$; busca, inserção no fim e remoção na lista encadeada são $O(n)$ sem ponteiro para o último nó.

\section*{Conclusão}
O artefato foi organizado para demonstrar tanto o conceito estudado quanto sua aplicação prática. As evidências de execução só são consideradas válidas quando produzidas pela ferramenta correspondente.
\end{document}
